\section{The first spectral support restriction}
\label{sec:fourier-support}
The joint boundary theorem also controls the Fourier support of the
continued distributions. We use the radian convention
\[
 \widehat f(X)=\int e^{-iX\cdot p}f(p)\,dp,\qquad
 \widehat W(f)=W(\widehat f).
\]
Thus the inverse test transform has factor $(2\pi)^{-4n}$. The dot product
here is the Euclidean pairing of the absolute coordinates, including each
real time coordinate; no relative-coordinate reduction has been made.
For a height tuple $y$, write $d(y)=(y_1e_0,\ldots,y_ne_0)$.

\begin{lemma}[Exponential testing identity]\label{lem:exponential-testing}
For strictly ordered $y$ and every smooth compactly supported frequency test
$f$,
\begin{equation}\label{eq:frequency-height}
 \widehat W_n(v)(f)
   =\psi_y\big(\widehat{\,e^{d(y)\cdot p}f(p)\,}\big).
\end{equation}
\end{lemma}
\begin{proof}
Set $f_s(p)=e^{s d(y)\cdot p}f(p)$ and $g_s=\widehat f_s$.
All $f_s$ have the same compact support. Differentiating their explicit
formula, including all spatial derivatives, shows that
$s\mapsto f_s$ is differentiable in the Schwartz topology with derivative
$(d(y)\cdot p)f_s$. To see the topology directly, every weighted derivative
is supported in that fixed compact set; Taylor's formula in $s$, uniformly
on it for each derivative, proves convergence in each Schwartz seminorm.
Continuity of the Fourier transform then gives
\[
 g_s'=\widehat{(d(y)\cdot p)f_s}=-D_yg_s.
\]
The sign follows by differentiating the kernel $e^{-iX\cdot p}$ and using
$D_y=-i\partial_y$.

For $s>0$, differentiate $\psi_{sy}(g_s)$. The varying-test step is
legitimate with the uniform seminorm bound proved in
Theorem~\ref{thm:joint-boundary}: write a difference quotient as the sum
of the change in the test at the new height and the change in height with
the old test fixed. The first part converges to $\psi_{sy}(g_s')$, because
the test quotient converges in Schwartz space and the family of functionals
has a common continuous seminorm bound. Lemma~\ref{lem:tested-cr} handles
the second part. Their sum is
$\psi_{sy}(-D_yg_s)+\psi_{sy}(D_yg_s)=0$.
The function is therefore constant on $(0,\infty)$. As $s\downarrow0$,
$g_s\to\widehat f$ in Schwartz space, and the same uniform bound plus the
joint boundary limit gives
$\psi_{sy}(g_s)\to W_n(v)(\widehat f)$. Comparing this with its value
at $s=1$ proves the formula.
\end{proof}

\begin{proposition}\label{prop:half-support}
For every label tuple, the distribution $\widehat W_n(v)$ is supported in
\[
 H_n=\left\{p:\sum_{i=1}^n y_i p_i^0\geq0
                    \text{ for every }y_1\leq\cdots\leq y_n\right\}.
\]
Equivalently,
\[
 \sum_{i=1}^n p_i^0=0,\qquad
 \sum_{i=j}^n p_i^0\geq0\quad(2\leq j\leq n).
\]
\end{proposition}
\begin{proof}
Fix a weakly ordered $y$ and a smooth compactly supported $f$ whose support
lies in the open set $d(y)\cdot p<-1$. Put
$y(t)=ty+r$, where $r=(0,1,\ldots,n-1)$ and $t\geq0$.
These heights are strictly ordered, with all gaps at least one. The
functions
\[
 h_t(p)=e^{d(y(t))\cdot p}f(p)
\]
tend to zero in Schwartz space. Indeed, on the fixed compact support every
derivative is a finite sum of a fixed smooth coefficient times a power of
$t$ times $e^{t d(y)\cdot p}$, and the latter is bounded by $e^{-t}$.
All polynomially weighted derivative suprema therefore tend to zero.
Equation~\eqref{eq:frequency-height} and the gap bound at $k=1$ imply
\[
 |\widehat W_n(v)(f)|
  =|\psi_{y(t)}(\widehat h_t)|
  \leq C\|\widehat h_t\|_{L^1}\longrightarrow0.
\]
Thus the distribution vanishes on that open set. If a point has
$d(y)\cdot p<0$, rescale $y$ by a positive constant so that its value is
less than $-1$. A neighborhood of the point then has the same strict
inequality. This proves support in $H_n$.

To identify $H_n$, the constant height tuples $1$ and $-1$ give zero total
time momentum. The weakly increasing step tuple equal to zero before $j$
and one from $j$ onward gives nonnegativity of each suffix sum. Conversely,
for any weakly ordered $y$,
\[
 \sum_i y_i p_i^0
 =y_1\sum_i p_i^0+
       \sum_{j=2}^n(y_j-y_{j-1})\sum_{i=j}^n p_i^0\geq0.
\]
This proves the equivalence, with the empty and single-slot conventions
included.
\end{proof}

\begin{corollary}[Conditional Lorentz improvement]\label{cor:cone-support}
If the boundary distributions are covariant under the continuous
finite-dimensional label representation of the proper orthochronous Lorentz
group, then
the support of every $\widehat W_n(v)$ lies in
\[
 \left\{p:\sum_i p_i=0,\quad \sum_{i=j}^n p_i\in\overline V_+
                                     \ (2\leq j\leq n)\right\},
 \qquad \overline V_+=\{p:p^0\geq|\mathbf p|\}.
\]
\end{corollary}
\begin{proof}
Fourier transformation of the covariance identity changes momenta by the
dual Lorentz action: this follows by the linear change of variables in the
test transform, whose Jacobian has absolute value one. The transformed
labels are again an allowed tuple, so Proposition~\ref{prop:half-support}
applies after every such dual Lorentz transformation. The inverse-transpose
action ranges over the same group: for the Minkowski metric matrix $\eta$,
$\Lambda^{-T}=\eta\Lambda\eta$, which preserves the future cone, and
this operation is an involution. In particular every transformed
suffix momentum has nonnegative time coordinate.

A real vector $q$ has this property exactly when $q\in\overline V_+$.
One implication follows from preservation of the future cone. For the
other, if $q^0<0$ the identity transformation suffices. If
$0\leq q^0<|\mathbf q|$, choose a boost in the direction of
$\mathbf q$ with speed $\beta$ strictly between $q^0/|\mathbf q|$ and
one; the transformed time coordinate
$(q^0-\beta|\mathbf q|)/\sqrt{1-\beta^2}$ is negative, a contradiction.
This proves the suffix assertion. The total momentum has zero time
coordinate in every frame. The identity frame makes its time coordinate
zero, and the same boost argument forces its spatial part to vanish.
\end{proof}

The Lorentz covariance hypothesis in this last corollary is still to be
proved from the infinitesimal Euclidean rotation identities. The conclusion
here is support of the distributions in momentum variables; constructing
the projection-valued momentum measure of the Hilbert-space representation
is a separate step.
